\subsection{Growing-dimensional support recovery}

This subsection proves the support identity for every growing-dimensional
regime in Theorem~\ref{thm:main-support-law}. In the Gaussian-bump coordinates
of Subsection~\ref{sec:mv-one-atom-geometry}, the buffered reductions below
identify every noncentral fitted atom with one radial exceedance while
retaining the exact centering correction.

\subsubsection{Exact edge coordinates}

Recall the exact radial coordinate \(\xi_{n,i}^\star\), scale
\(\delta_n^\star\), and identity \eqref{eq:exact-support-radius} from
Section~\ref{sec:mv-gamma-normalization}.
For fixed \(C,J<\infty\) and \(\eta>0\), recall the buffered event
\(\mathcal B_n(C,J,\eta)\) from \eqref{eq:exact-support-buffer}.
The centered radial Poisson limit makes the complement of such buffered
events negligible after \(C,J\to\infty\) and then \(\eta\downarrow0\).
It therefore remains to prove that every positive buffered excess creates
exactly one noncentral support point.

When \(d_nL_n=o(n^2)\), the exact coordinate agrees with the
\(2\log n\)-based coordinate used in Sections~\ref{sec:radial-transition}
and~\ref{sec:proof-program}. The next
elementary comparison lets us invoke the radial estimates from those sections
without carrying two normalizations through each proof.

\begin{lemma}
\label{lem:exact-old-edge-coordinate}
Suppose \(d_nL_n=o(n^2)\) and \(s_n^\star=O(1)\).  With \(s_n\),
\(\xi_{n,i}\), and \(\delta_n=a_n/\gamma_n\) as in
\eqref{eq:s-definition} and \eqref{eq:bump-height-xi},
\[
 s_n-s_n^\star=o(1),
 \qquad
 \frac{\delta_n^\star}{\delta_n}=1+o(1).
\]
Moreover, for every fixed \(C<\infty\),
\begin{equation}
\label{eq:exact-old-edge-comparison}
 \max_{i:\,|\xi_{n,i}^\star|\le C}
 \left|\xi_{n,i}^\star-(\xi_{n,i}-s_n)\right|
 \longrightarrow_{\Pbb}0.
\end{equation}
\end{lemma}

\begin{proof}
The second assertion is immediate from
\(\delta_n^\star=\sigma_n^2a_n/\gamma_n\).  The definitions give
\[
 s_n-s_n^\star
 =\frac{\gamma_n}{a_n}
   \left(L_n-\frac{R_n^\star}{\sigma_n^2}\right).
\]
Since \(L_n-R_n^\star=O(n^{-1})\) and
\(1-\sigma_n^2=n^{-1}\), its absolute value is
\(O\{b_n/(na_n)+\gamma_n/(na_n)\}\).  This tends to zero: it is
\(O(L_n/n)\) when \(d_n=O(L_n)\), while the superlogarithmic Gamma
expansions make it \(O(\sqrt{d_nL_n}/n)=o(1)\) otherwise.

In fact, the two excess coordinates satisfy the exact identity
\begin{equation}
\label{eq:exact-old-excess-identity}
 \xi_{n,i}-s_n
 =\sigma_n^2\xi_{n,i}^\star
  -\frac{L_n-R_n^\star}{\delta_n}.
\end{equation}
Now \(L_n-R_n^\star=O(n^{-1})\).  Also
\((n\delta_n)^{-1}=o(1)\): this is immediate when \(d_n=O(L_n)\), and in
the superlogarithmic range it is
\(O\{\sqrt{d_n/L_n}/n\}=o(L_n^{-1})\).
Equation \eqref{eq:exact-old-excess-identity}, together with
\(\sigma_n^2=1-n^{-1}\), proves
\eqref{eq:exact-old-edge-comparison}.
\end{proof}


For later use, if \(c_1,\ldots,c_n>0\), define
\[
 \Gamma_c(w)=\frac1n\sum_{i=1}^n
 c_i\exp\{w\cdot q_i-|w|^2/2\}.
\]
The Gibbs variational formula gives the weighted entropy dual
\begin{equation}
\label{eq:weighted-entropy-dual}
 \log\Gamma_c(w)
 =\sup_{p\in\Delta_n}
 \left\{
  w\cdot\sum_ip_iq_i-\frac{|w|^2}{2}
  -K(p\Vert\mathbf u_n)+\sum_ip_i\log c_i
 \right\}.
\end{equation}

The entropy argument used for one-atomicity must be made stable under the
fitted denominators before it can count contacts.  The following weighted
version supplies precisely that stability.  It is stated only in the random
entropy regime, where it will be used.

\begin{lemma}
\label{lem:weighted-entropy-localization}
Suppose \eqref{eq:entropy-regime} holds, \(s_n=O(1)\), and
\(d_nL_n=o(n^2)\).  Put
\[
 q_i=\frac{Z_i-\bar Z_n}{\sqrt{\gamma_n}},
 \qquad
 R_i=\frac{|q_i|^2}{2},
 \qquad
 \delta_n=\frac{a_n}{\gamma_n},
 \qquad
 \mu_n^{\rm off}=\max_{i\ne j}|q_i\cdot q_j|,
 \qquad
 \mathfrak k_n=\frac{\|(q_i\cdot q_j)_{i,j\le n}\|_{\rm op}}{2n}.
\]
Fix \(C,J<\infty\) and \(\eta>0\), and suppose
\[
 \max_i(R_i-L_n)/\delta_n\le C,\qquad
 \#\{i:R_i>L_n-\eta\delta_n\}\le J,\qquad
 \min_i|R_i-L_n|>\eta\delta_n.
\]
Let \(A=\{i:R_i>L_n\}\).  Suppose a random vector
\(\beta=(\beta_1,\ldots,\beta_n)\) satisfies, uniformly on this event,
\begin{equation}
\label{eq:weighted-entropy-alpha-profile}
 \beta_i=-(R_i-L_n)+e_{n,i}\quad(i\in A),
 \qquad
 \max_{i\in A}|e_{n,i}|
 +\max_{j\notin A}|\beta_j|
 =:\omega_n,
\end{equation}
where
\begin{equation}
\label{eq:weighted-entropy-omega-rate}
 \omega_n
 =
 O_{\Pbb}\left\{\frac{e^{\mu_n^{\rm off}}}n+\frac{\delta_nL_n}{n}\right\}
 =n^{-1+o_{\Pbb}(1)}
 =o_{\Pbb}(\delta_n).
\end{equation}
If
\[
 F_\beta(p)
 =
 \frac12\left|\sum_i p_iq_i\right|^2
 -K(p\Vert \mathbf u_n)+\sum_i p_i\beta_i,
\]
then, with probability tending to one, every \(p\in\Delta_n\) satisfying
\(F_\beta(p)\ge0\) has one of the following two forms:
\begin{equation}
\label{eq:weighted-entropy-two-patches}
 \|p-\mathbf u_n\|_1\le n^{-1/2+o_{\Pbb}(1)},
 \qquad\text{or}\qquad
 1-p_j\le n^{-1/2+o_{\Pbb}(1)}
 \quad\text{for some }j\in A.
\end{equation}
In the first case \(|\sum_i p_iq_i|=o_{\Pbb}(1)\); in the second,
\(|\sum_i p_iq_i-q_j|=o_{\Pbb}(1)\). Moreover, if \(p=p(w)\) is the
Gibbs vector of a weighted certificate with value at least one, then \(w\)
lies within \(n^{-1/2+o_{\Pbb}(1)}\) of zero in the first case and of
\(q_j\) in the second.
\end{lemma}

\begin{proof}
Lemma~\ref{lem:random-off-diagonal} and \eqref{eq:entropy-regime} imply that
the deterministic bounds in Proposition~\ref{prop:refined-deterministic-gram}
hold with probability tending to one. Replace
\(\mu_n^{\rm off}\) and \(\upsilon_n\) by deterministic upper bounds that
hold with probability tending to one and preserve their asymptotic
separation. For readability, continue to denote these deterministic bounds by
\(\mu_n^{\rm off}\) and \(\upsilon_n\). We may then choose
deterministic \(a_n^\circ\) so that
\[
 1+\mu_n^{\rm off}\ll a_n^\circ\ll
 \min\{\log(1/\upsilon_n),\log(n\delta_n),L_n\},
\]
on an event of probability tending to one, and then set
\[
 T_n=e^{a_n^\circ},
 \qquad
 L_ny_n=\{a_n^\circ\log(n\delta_n)\}^{1/2}.
\]
These choices satisfy \eqref{eq:refined-gram-conditions}. Since
\(\max_iR_i=L_n+O(\delta_n)\), the extra diagonal term in
\eqref{eq:heavy-energy} is
\[
 O(\delta_n)\sum_i b_i^2
 =o\left(L_n\sum_i b_i^2\right)
\]
and is absorbed by the heavy-entropy term in
\eqref{eq:heavy-entropy-scalar}.

First suppose \(\max_i p_i\le1-y_n\).  Repeat the light-heavy
decomposition in Proposition~\ref{prop:refined-deterministic-gram}. Explicitly,
put
\[
 f_i=np_i,\qquad \ell_i=f_i\wedge T_n,\qquad h_i=(f_i-T_n)_+,
 \qquad b_i=\frac{h_i}{n},
 \qquad K_0=\frac1n\sum_i\phi(\ell_i),\qquad
 z=\sum_i b_i,
\]
where \(\phi(x)=x\log x-x+1\).
The active diagonal excess is canceled by the active fitted weight:
because \(b_i^2\le p_i\),
\[
 \sum_{i\in A}(R_i-L_n)b_i^2+\sum_i p_i\beta_i
 \le\omega_n.
\]
Keeping half of the slack in \eqref{eq:light-energy} and
\eqref{eq:heavy-entropy-scalar} therefore gives
\begin{equation}
\label{eq:weighted-diffuse-slack}
 F_\beta(p)
 \le
 -cK_0-ca_n^\circ z+\omega_n.
\end{equation}
Thus \(F_\beta(p)\ge0\) implies
\(K_0+a_n^\circ z=O_{\Pbb}(\omega_n)\).  Moreover,
\[
 \|p-\mathbf u_n\|_1
 \le
 C\left\{
   \left(\frac{T_n\omega_n}{a_n^\circ}\right)^{1/2}
   +\frac{\omega_n}{a_n^\circ}
 \right\}
 \le n^{-1/2+o_{\Pbb}(1)}.
\]
The last inequality uses \(T_n=n^{o(1)}\) and
\eqref{eq:weighted-entropy-omega-rate}; in particular,
\(L_n(T_n\omega_n/a_n^\circ)^{1/2}=o_{\Pbb}(1)\).
The Gram bound then gives
\(|\sum_i p_iq_i|=o_{\Pbb}(1)\), proving the first alternative.

It remains to consider \(p_i=1-y>1-y_n\).  Write
\(\sum_jp_jq_j=(1-y)q_i+y\bar q\), as in the last part of the proof of
Proposition~\ref{prop:refined-deterministic-gram}.  Put
\[
 \mathfrak H_n
 =
 L_n-\mu_n^{\rm off}-O(\delta_n)-2L_ny_n-\omega_n.
\]
The preceding choices give
\(\mathfrak H_n\sim L_n\) and
\(e^{-\mathfrak H_n}=n^{-1+o(1)}=o(\delta_n)\). Convexity,
the entropy bound \(h_{\rm bin}(y)\le y\log(e/y)\), and the fitted weights give
\begin{equation}
\label{eq:weighted-near-vertex}
 F_\beta(p)
 \le
 \alpha_i
 +y\left\{\log\frac ey-\mathfrak H_n\right\},
\end{equation}
where \(\alpha_i=R_i-L_n+\beta_i\).

If \(i\notin A\), the first term on the right of
\eqref{eq:weighted-near-vertex} is at most
\(-\eta\delta_n+\omega_n\). The second term on the right of
\eqref{eq:weighted-near-vertex} is at most
\(e^{-\mathfrak H_n}=o(\delta_n)\).
Hence \(F_\beta(p)<0\).

If \(i\in A\), then \(|\alpha_i|\le\omega_n\).  For
\(y>e^{-\mathfrak H_n/2}\),
\[
 \log(e/y)-\mathfrak H_n\le-\mathfrak H_n/3
\]
for all large \(n\).  Nonnegative free energy consequently forces
\[
 y\le e^{-\mathfrak H_n/2}+\frac{3\omega_n}{\mathfrak H_n}
 \le n^{-1/2+o_{\Pbb}(1)}.
\]
Since \(\max_j|q_j|=O_{\Pbb}(\sqrt{L_n})\), this gives
\[
 \left|\sum_jp_jq_j-q_i\right|
 =y|\bar q-q_i|=o_{\Pbb}(1),
\]
which is the second alternative in
\eqref{eq:weighted-entropy-two-patches}.

For later use, define the weighted certificate and its Gibbs vector by
\[
 \Gamma_\beta(w)
 =\frac1n\sum_{i=1}^n
   \exp\left\{w\cdot q_i-\frac{|w|^2}{2}+\beta_i\right\},
 \qquad
 p_i(w)
 =\frac{\exp\{w\cdot q_i+\beta_i\}}
        {\sum_j\exp\{w\cdot q_j+\beta_j\}}.
\]
The exact identity
\[
 F_\beta\{p(w)\}-\log\Gamma_\beta(w)
 =
 \frac12\left|w-\sum_i p_i(w)q_i\right|^2
\]
shows that the two alternatives in
\eqref{eq:weighted-entropy-two-patches} also localize \(w\) whenever
\(\Gamma_\beta(w)\ge1\). On the diffuse branch,
\eqref{eq:weighted-diffuse-slack} gives
\(F_\beta\{p(w)\}=n^{-1+o_{\Pbb}(1)}\), while the Gram estimate and
\eqref{eq:weighted-entropy-two-patches} give
\(\lvert\sum_i p_i(w)q_i\rvert=n^{-1/2+o_{\Pbb}(1)}\). On an active
branch, \eqref{eq:weighted-near-vertex} gives
\(F_\beta\{p(w)\}=n^{-1+o_{\Pbb}(1)}\), and
\(\lvert\sum_i p_i(w)q_i-q_j\rvert=n^{-1/2+o_{\Pbb}(1)}\). The displayed
identity therefore proves the asserted spatial localization.
Moreover,
\[
 \nabla^2\log\Gamma_\beta(w)
 =\operatorname{Cov}_{p(w)}(q_i)-I.
\]
On the diffuse branch the covariance norm is at most
\(2\mathfrak k_n+O\{L_n\|p-\mathbf u_n\|_1\}=o_{\Pbb}(1)\); on an active branch it is
\(O\{L_n(1-p_i)\}=o_{\Pbb}(1)\).  Thus every relevant patch is strictly
log-concave for all large \(n\).
\end{proof}

The slowly growing range needs a local version of the truncated shell
argument.  Conditional first moments are enough.  Exponential tilting
shows that the supremum of the background on a unit edge patch has mean
\(\exp\{-\eps_nL_n+O(\sqrt{d_n}+\eps_n\sqrt{L_n})\}\), which beats the
number of upper centers even when \(d_n\) grows arbitrarily slowly.

\begin{lemma}
\label{lem:slow-punctured-c3}
Suppose
\[
 d_n\longrightarrow\infty,\qquad
 d_n\le(\log L_n)^2,\qquad s_n=O(1),
\]
and put \(A_n=\sqrt{d_n}\).  There is a random set
\(\mathcal U_n\), containing every index with
\(\xi_{n,i}-s_n>-\eta\), such that, with probability tending to one,
\[
 |\mathcal U_n|\le e^{2A_n},\qquad
 \min_{\substack{i,j\in\mathcal U_n\\i\ne j}}
 |q_i-q_j|^2=4L_n\{1+o(1)\},
\]
and, for every fixed \(R<\infty\) and \(0\le k\le3\),
\begin{equation}
\label{eq:slow-punctured-c3}
 \max_{j\in\mathcal U_n}\ \sup_{|w-q_j|\le R}
 \left\|
  \nabla^k\sum_{i\notin\mathcal U_n}
  h_ie^{-|w-q_i|^2/2}
 \right\|=o_{\Pbb}(1).
\end{equation}
\end{lemma}

\begin{proof}
Work first with the independent uncentered vectors
\(Q_i=Z_i/\sqrt{\gamma_n}\), and let \(\mathcal U_n^\circ\) contain the
indices whose uncentered radial coordinate is at least \(-A_n-1\).
Write \(h_i^\circ=n^{-1}e^{|Q_i|^2/2}\) for the corresponding uncentered
bump height.
The Gamma hazard estimate and Markov's inequality give
\[
 |\mathcal U_n^\circ|\le e^{2A_n}
\]
with probability tending to one. Conditional on the radii, their
directions are independent and uniform. Since
\(\log|\mathcal U_n^\circ|=O(A_n)=o(d_n)\), spherical concentration gives
the displayed separation, first for the \(Q_i\)'s and then for the
centered \(q_i\)'s.

Conditional on \(Q_j\), \(j\in\mathcal U_n^\circ\), all remaining vectors
are independent of \(Q_j\).  Write \(w=Q_j+h\), \(|h|\le R\).  For every
multi-index \(|\alpha|\le3\), exponential tilting from
\(\mathcal N(0,\gamma_n^{-1}I)\) to
\(\mathcal N(Q_j/\gamma_n,\gamma_n^{-1}I)\) gives
\begin{align*}
 &\sum_{i\ne j}\E\left[
  \sup_{|h|\le R}
  \left|\partial^\alpha
  \{h_i^\circ e^{-|Q_j+h-Q_i|^2/2}\}\right|
  \mathbf1_{\{i\notin\mathcal U_n^\circ\}}
  \,\middle|\,Q_j\right]\\
 &\qquad\le
 C_{\alpha,R}
 \exp\left\{
  -\frac{\eps_n|Q_j|^2}{2}
  +C_{\alpha,R}\bigl(\sqrt{d_n}+\eps_n|Q_j|+1\bigr)
 \right\}.
\end{align*}
Indeed, after tilting,
\(Q_i-Q_j\) is Gaussian with mean \(-\eps_nQ_j\), covariance
\(\gamma_n^{-1}I\), and the supremum over \(h\) contributes at most a
fixed polynomial in its norm times \(e^{R|Q_i-Q_j|}\).
Uniformly over \(j\in\mathcal U_n^\circ\),
\(|Q_j|^2/2=L_n+O(A_n)\), while
\[
 \eps_nL_n=(1-o(1))\frac{d_n}{2}\log\frac{2L_n}{d_n}.
\]
This dominates
\(A_n+\sqrt{d_n}+\eps_n\sqrt{L_n}\), since
\(d_n\to\infty\) and \(d_n\le(\log L_n)^2\).  Conditional Markov's
inequality is applied to the random set through
\[
 \Pbb\{\exists j\in\mathcal U_n^\circ:B_{n,j}\}
 \le
 \sum_{j=1}^n
 \E\left[
  \mathbf1_{\{j\in\mathcal U_n^\circ\}}
  \Pbb(B_{n,j}\mid Q_j)
 \right],
\]
where \(B_{n,j}\) is the failure of the displayed derivative bound.
Since
\(\E|\mathcal U_n^\circ|=e^{(1+o(1))A_n}\), the conditional exponential
moment bound in the preceding display makes this sum \(o(1)\).  Passing from coordinate derivatives
to the tensor norm costs at most \(O(d_n^k)\) for order \(k\le3\), whose
logarithm is negligible compared with \(\eps_nL_n\).  This proves the
uncentered version of \eqref{eq:slow-punctured-c3}.

Finally,
\[
 B_n(w)=e^{-w\cdot\bar Z_n/\sqrt{\gamma_n}}B_n^\circ(w),
 \qquad
 \frac{|\bar Z_n|}{\sqrt{\gamma_n}}
 =O_{\Pbb}\{\sqrt{d_n/n}\}.
\]
On the relevant balls \(|w|=O(\sqrt{L_n})\), this multiplicative factor
is \(1+o_{\Pbb}(1)\), while each of its first four derivatives is
\(o_{\Pbb}(1)\).  The uniform
centered--uncentered radial difference is \(o_{\Pbb}(1)\), so the
one-unit cutoff buffer shows that
\(\mathcal U_n:=\mathcal U_n^\circ\) contains every centered index above
\(-\eta\).  This proves all three assertions.
\end{proof}

\subsubsection{Uniform finite-contact continuation}

The support-count argument needs a joint expansion over all active edge
patches.  The following lemma supplies it.  The hypotheses are stated in
terms of absolute derivative envelopes; this is important because the
reciprocal fitted values used in the KKT certificate need not preserve
cancellation among derivatives of different bumps.

\begin{lemma}
\label{lem:uniform-finite-edge-expansion}
Let \(L=\log n\), and let \(q_1,\ldots,q_n\in\mathbb R^{d_n}\) satisfy
\[
 \sum_{j=1}^nq_j=0,\qquad \max_j|q_j|\le C_0\sqrt L.
\]
Suppose
\[
 \frac{|q_i|^2}{2}=L+\delta v_i,\qquad 0<\delta\le2,
 \qquad \bar\delta=1\wedge\delta.
\]
Let \(\mathcal U\subseteq[n]\) contain every \(i\) with \(v_i>-\eta\),
put \(A=\{i:v_i>0\}\), and assume
\[
 |A|\le J,\qquad \eta\le v_i\le C\quad(i\in A),\qquad
 v_i\le-\eta\quad(i\in\mathcal U\setminus A).
\]
For \(0\le k\le3\), define
\begin{equation}
\label{eq:finite-edge-derivative-envelope}
\mathcal E_{i,k}(r)
 =
 \sup_{|z|\le r}\frac1n\sum_{j\ne i}
 \{1+|q_j-q_i-z|^k\}
 \exp\left\{(q_i+z)\cdot q_j-\frac{|q_i+z|^2}{2}\right\}.
\end{equation}
Fix \(0<r_n\le R_n\le1/4\), and put
\(e_{k,n}=\max_{i\in\mathcal U}\mathcal E_{i,k}(2R_n)\).  Suppose that
\begin{equation}
\label{eq:finite-edge-envelope-assumptions}
 \max_{i\in\mathcal U}\mathcal E_{i,0}(2R_n)
   =o_{\Pbb}(\bar\delta),\qquad
 \max_{i\in\mathcal U}\mathcal E_{i,1}(2R_n)
   =o_{\Pbb}(r_n),
\end{equation}
and
\begin{equation}
\label{eq:finite-edge-higher-envelope-assumptions}
 \begin{aligned}
 \max_{i\in\mathcal U}
 \{\mathcal E_{i,2}(2R_n)+\mathcal E_{i,3}(2R_n)\}
 &=o_{\Pbb}(1),\qquad
 n\bar\delta\longrightarrow\infty,
 \qquad
 r_n^2=o(\bar\delta),\\
 \frac{\bar\delta L}{n}&=o(r_n),\qquad
 r_ne_{1,n}+r_n^2e_{2,n}+r_n^3e_{3,n}
 &=o_{\Pbb}(\bar\delta).
 \end{aligned}
\end{equation}

For \(\tau\in\mathbb R^{d_n}\), \(z_i\in\mathbb R^{d_n}\), and
\(m_i\ge0\), set
\begin{equation}
\label{eq:finite-edge-candidate}
 \Pi_{\tau,z,m}
 =
 \left(1-\frac{\sum_{i\in A}m_i}{n}\right)\delta_\tau
 +\frac1n\sum_{i\in A}m_i\delta_{q_i+z_i}.
\end{equation}
Uniformly for
\[
 |\tau|\le K\bar\delta\sqrt L/n,\qquad |z_i|\le R_n,\qquad
 0\le m_i\le K\bar\delta,
\]
its relative log likelihood has the expansion
\begin{equation}
\label{eq:joint-finite-edge-expansion}
 \ell_q(\Pi_{\tau,z,m})
 =
 -\frac n2|\tau|^2
 +\sum_{i\in A}
 \left[
  \log\{1+m_i e^{\vartheta_i(\tau,z_i)}\}-m_i
 \right]
 +\mathcal R_n(\tau,z,m),
\end{equation}
where
\[
 \vartheta_i(\tau,z)
 =
 \delta v_i-\tau\cdot q_i+\frac{|\tau|^2}{2}-\frac{|z|^2}{2}.
\]
The remainder is \(o_{\Pbb}(\bar\delta^2)\), uniformly on this set.  The
same expansion may be differentiated twice in the mass variables and
three times in the spatial variables.  In particular, uniformly,
\begin{align}
\label{eq:joint-finite-edge-derivatives}
 \|\partial_m\mathcal R_n\|
   &\le O_{\Pbb}(e_{0,n}+n^{-1})=o_{\Pbb}(\bar\delta),
 &\|\partial_m^2\mathcal R_n\|&=o_{\Pbb}(1),\notag\\
 \|\nabla_z\mathcal R_n\|
   &\le O_{\Pbb}(\bar\delta e_{1,n})
     =o_{\Pbb}(\bar\delta r_n),
 &\|\partial_m\nabla_z\mathcal R_n\|
   &=O_{\Pbb}(e_{1,n})=o_{\Pbb}(r_n),\notag\\
 \|\nabla_z^2\mathcal R_n\|+
 \|\nabla_z^3\mathcal R_n\|
   &\le O_{\Pbb}\{\bar\delta(e_{2,n}+e_{3,n})\}
     =o_{\Pbb}(\bar\delta),
 &\|\nabla_\tau\mathcal R_n\|
   &=o_{\Pbb}(\bar\delta\sqrt L),
\end{align}
with all mixed blocks interpreted in the maximum block norm.

There is a critical point of \(\ell_q\) in these patches, unique near the
locations and masses displayed in
\eqref{eq:finite-edge-contact-asymptotics}--\eqref{eq:finite-edge-mass-asymptotics},
for which
\begin{equation}
\label{eq:finite-edge-contact-asymptotics}
 \tau_n
 =
 -\frac1n\sum_{i\in A}(1-e^{-\delta v_i})q_i
 +o_{\Pbb}(\bar\delta\sqrt L/n),\qquad
 z_{n,i}=o_{\Pbb}(r_n),
\end{equation}
and
\begin{equation}
\label{eq:finite-edge-mass-asymptotics}
 m_{n,i}=1-e^{-\delta v_i}+o_{\Pbb}(\bar\delta).
\end{equation}
In particular, all edge masses are positive.  If
\(\widetilde\Pi=\Pi_{\tau_n,z_n,m_n}\), then
\[
 \Gamma_{\widetilde\Pi}(\tau_n)
 =\Gamma_{\widetilde\Pi}(q_i+z_{n,i})=1,\qquad
 \nabla\Gamma_{\widetilde\Pi}(\tau_n)
 =\nabla\Gamma_{\widetilde\Pi}(q_i+z_{n,i})=0.
\]
Writing \(c_j=M_j(\widetilde\Pi)^{-1}\), one has
\begin{equation}
\label{eq:finite-edge-denominator-profile}
 \max_{j\notin A}(\log c_j)_+
 =O_{\Pbb}(\bar\delta L/n),\qquad
 \log c_i=-\delta v_i+o_{\Pbb}(\bar\delta)\quad(i\in A).
\end{equation}
Moreover, each active patch has the \(C^2\) representation
\begin{equation}
\label{eq:finite-edge-certificate-c2}
 \Gamma_{\widetilde\Pi}(q_i+y)
 =
 A_{n,i}e^{-|y|^2/2}+E_{n,i}(y),\qquad
 A_{n,i}=1+o_{\Pbb}(1),
\end{equation}
where
\[
 \|E_{n,i}\|_{C^0(B(0,R_n))}=o_{\Pbb}(\bar\delta),
 \qquad
 \|E_{n,i}\|_{C^2(B(0,R_n))}=o_{\Pbb}(1).
\]
Thus the certificate is strictly log-concave on every active patch and
has exactly one level-one point there.  On every inactive patch indexed
by \(\mathcal U\setminus A\),
\[
 \sup_{|y|\le R_n}\Gamma_{\widetilde\Pi}(q_i+y)
 \le1-c_\eta\bar\delta.
\]
\end{lemma}

\begin{proof}
Write \(k_w(j)=\exp\{w\cdot q_j-|w|^2/2\}\), \(m_+=\sum_{i\in A}m_i\),
and \(w_i=q_i+z_i\).  Factoring out the moving central component gives
the exact identity
\[
 M_j(\Pi_{\tau,z,m})=k_\tau(j)(1+x_j),\qquad
 x_j=-\frac{m_+}{n}
 +\frac1n\sum_{i\in A}m_iK_{ij},
\]
where
\[
 K_{ij}
 =\frac{k_{w_i}(j)}{k_\tau(j)}
 =\exp\left\{(w_i-\tau)\cdot(q_j-\tau)
              -\frac{|w_i-\tau|^2}{2}\right\}.
\]
At the observation belonging to the \(i\)th patch,
\begin{equation}
\label{eq:finite-edge-own-response}
 \frac{K_{ii}}n=e^{\vartheta_i(\tau,z_i)}.
\end{equation}
For \(j\notin A\), positivity and
\eqref{eq:finite-edge-envelope-assumptions} give, uniformly over the
displayed parameter set,
\[
 \max_{j\notin A}|x_j|=o_{\Pbb}(1),\qquad
 \sum_{j\notin A}x_j=-m_++o_{\Pbb}(\bar\delta^2),
 \qquad
 \sum_{j\notin A}x_j^2=o_{\Pbb}(\bar\delta^2).
\]
Indeed, after putting \(a_i=w_i-\tau\), the positive response away from
the own observation equals
\[
 \frac{m_i}{n}\sum_{j\ne i}K_{ij}
 =
 m_i e^{-a_i\cdot\tau}
 \frac1n\sum_{j\ne i}e^{a_i\cdot q_j-|a_i|^2/2}
 =o_{\Pbb}(\bar\delta^2).
\]
The negative part is \(m_+/n\) per observation.  Omitting the at most
\(J\) active observations therefore costs \(O(\bar\delta/n)\), which is
\(o(\bar\delta^2)\) by \(n\bar\delta\to\infty\).  Positivity also gives
\[
 \frac1{n^2}\sum_{j\ne i}K_{ij}^2
 \le
 \left(\frac1n\max_{j\ne i}K_{ij}\right)
 \left(\frac1n\sum_{j\ne i}K_{ij}\right)
 \le O_{\Pbb}(e_{0,n}^2),
\]
which controls the quadratic Taylor remainder. Applying the same product
bound after inserting the polynomial factors from
\eqref{eq:finite-edge-derivative-envelope} yields
\begin{align}
 |\mathcal R_n|
 &\le O_{\Pbb}\{\bar\delta e_{0,n}+e_{0,n}^2
                   +\bar\delta/n\},\notag\\
 \|\partial_m\mathcal R_n\|
 &\le O_{\Pbb}(e_{0,n}+n^{-1}),
 &\|\partial_m^2\mathcal R_n\|&=o_{\Pbb}(1),\notag\\
 \|\nabla_z\mathcal R_n\|
 &\le O_{\Pbb}(\bar\delta e_{1,n}),
 &\|\partial_m\nabla_z\mathcal R_n\|
   &\le O_{\Pbb}(e_{1,n}),\notag\\
 \|\nabla_z^2\mathcal R_n\|+
 \|\nabla_z^3\mathcal R_n\|
 &\le O_{\Pbb}\{\bar\delta(e_{2,n}+e_{3,n})\}.
 \label{eq:finite-edge-remainder-bounds}
\end{align}
Differentiation with respect to \(\tau\) produces the same relative
displacements, together with terms bounded by \(\sqrt L/n\) from the
finitely many omitted observations.  Hence
\begin{equation}
\label{eq:finite-edge-tau-remainder}
 \|\nabla_\tau\mathcal R_n\|
 \le O_{\Pbb}\{\bar\delta e_{1,n}
                +\bar\delta\sqrt L\,e_{0,n}
                +\bar\delta\sqrt L/n\}
 =o_{\Pbb}(\bar\delta\sqrt L).
\end{equation}
These estimates prove the remainder assertions in
\eqref{eq:joint-finite-edge-derivatives}.

Since \(\sum_jq_j=0\),
\(\sum_j\log k_\tau(j)=-n|\tau|^2/2\).  Taylor expansion of
\(\sum_{j\notin A}\log(1+x_j)\), together with
\eqref{eq:finite-edge-own-response}, proves
\eqref{eq:joint-finite-edge-expansion}.  Differentiating the same exact
factorization proves \eqref{eq:joint-finite-edge-derivatives}.  The
finiteness of \(|A|\) makes every bound uniform over simultaneous masses
and locations.

It remains to justify the critical point without treating the fitted
denominators as fixed.  Define the contact map
\[
 \mathcal F(\tau,z,m)
 =
 \left(
  \nabla\Gamma_\Pi(\tau),\
  \{\nabla\Gamma_\Pi(q_i+z_i)\}_{i\in A},\
  \{\Gamma_\Pi(q_i+z_i)-\Gamma_\Pi(\tau)\}_{i\in A}
 \right).
\]
Put \(m_i^0=1-e^{-\delta v_i}\),
\(\tau^0=-n^{-1}\sum_i m_i^0q_i\), and
\(\widetilde m_i^0=1-e^{-\vartheta_i(\tau^0,0)}\).  Then
\[
 |\tau^0|=O(\bar\delta\sqrt L/n),\qquad
 \widetilde m_i^0-m_i^0=o(\bar\delta),
\]
and \(m_i^0\ge c_{\eta,C}\bar\delta\).  Use the anisotropic norms
\[
 \|(f_0,\{f_i\},\{g_i\})\|_{\rm out}
 =
 \frac{n|f_0|}{\bar\delta\sqrt L}
 +\max_i\frac{|f_i|}{r_n}
 +\max_i\frac{|g_i|}{\bar\delta}
\]
and
\[
 \|(\Delta\tau,\Delta z,\Delta m)\|_{\rm in}
 =
 \frac{n|\Delta\tau|}{\bar\delta\sqrt L}
 +\max_i\frac{|\Delta z_i|}{r_n}
 +\max_i\frac{|\Delta m_i|}{\bar\delta}.
\]
The likelihood derivatives and contact equations are related by
\[
 \partial_\tau\ell_q
 =n\left(1-\frac{m_+}{n}\right)\nabla\Gamma_\Pi(\tau),
 \qquad
 \nabla_{z_i}\ell_q=m_i\nabla\Gamma_\Pi(q_i+z_i),
 \qquad
 \partial_{m_i}\ell_q
 =\Gamma_\Pi(q_i+z_i)-\Gamma_\Pi(\tau).
\]

For the quantitative estimates, put
\[
 \begin{split}
 \chi_n={}&
 \frac{e_{0,n}+n^{-1}}{\bar\delta}
 +\frac{e_{1,n}}{r_n}
 +\frac{e_{0,n}+e_{1,n}}{\sqrt L}
 +\frac Ln,\\
 \zeta_n={}&
 \frac{e_{0,n}}{\bar\delta}
 +\frac{e_{1,n}}{r_n}+e_{2,n}+e_{3,n}
 +\frac{r_ne_{1,n}+r_n^2e_{2,n}+r_n^3e_{3,n}}{\bar\delta}\\
 &+\frac1{n\bar\delta}+\frac{\bar\delta L}{nr_n}.
 \end{split}
\]
Every term is \(o_{\Pbb}(1)\) by the hypotheses.  At
\((\tau^0,0,\widetilde m^0)\), the principal part of
\eqref{eq:joint-finite-edge-expansion} is stationary in the spatial and
mass coordinates.  Its central derivative is
\(O(\bar\delta L^{3/2}/n)\).  Hence
\eqref{eq:finite-edge-remainder-bounds} and
\eqref{eq:finite-edge-tau-remainder} give the explicit residual estimate
\begin{equation}
\label{eq:finite-edge-scaled-residual}
 \|\mathcal F(\tau^0,0,\widetilde m^0)\|_{\rm out}
 =O_{\Pbb}(\chi_n)=o_{\Pbb}(1).
\end{equation}

Let \(\alpha_i=e^{-\vartheta_i(\tau^0,0)}\), which is bounded above and
below, and define
\[
 \mathcal J_n(\Delta\tau,\Delta z,\Delta m)
 =
 \left(
 -\Delta\tau-\frac1n\sum_i\alpha_iq_i\Delta m_i,\
 \{-\Delta z_i\}_{i\in A},\
 \{-\Delta m_i-\alpha_iq_i\cdot\Delta\tau\}_{i\in A}
 \right).
\]
Direct differentiation of the exact factorization gives
\begin{equation}
\label{eq:finite-edge-scaled-jacobian}
 \|D\mathcal F(\tau^0,0,\widetilde m^0)-\mathcal J_n\|_{{\rm in}\to{\rm out}}
 =O_{\Pbb}(\zeta_n)=o_{\Pbb}(1).
\end{equation}
The operator \(\mathcal J_n\) and its inverse are uniformly bounded.  Indeed,
its only nontrivial Schur correction is bounded by
\[
 \left\|\frac1n\sum_{i\in A}\alpha_i^2q_iq_i^\top\right\|_{\rm op}
 =O_J(L/n)=o(1).
\]

The own-response term requires one additional contribution in the
Jacobian-variation estimate.  On a scaled ball of radius \(\rho\), one has
\(|z_i|\le\rho r_n\), and differentiating
\(-|z_i|^2/2\) in the mass equation contributes at most
\(C\rho r_n^2/\bar\delta\) in the scaled operator norm.  The remaining
principal and remainder terms are controlled by \(\rho\) and \(\zeta_n\).
Consequently, uniformly for \(0<\rho\le1\),
\begin{equation}
\label{eq:finite-edge-scaled-jacobian-variation}
 \sup_{\|(\tau,z,m)-(\tau^0,0,\widetilde m^0)\|_{\rm in}\le\rho}
 \|D\mathcal F-D\mathcal F(\tau^0,0,\widetilde m^0)\|_{{\rm in}\to{\rm out}}
 \le
 O_{\Pbb}\left\{\zeta_n+\rho+
                  \frac{\rho r_n^2}{\bar\delta}\right\}.
\end{equation}
Since \(r_n^2=o(\bar\delta)\), the right side is
\(o_{\Pbb}(1)+O(\rho)\).

We finally verify rather than assume the existence of a vanishing Newton
radius.  Since \(\chi_n\xrightarrow{\Pbb}0\), choose integers \(N_k\uparrow\infty\)
so that
\[
 \Pbb\{\chi_n>k^{-2}\}\le k^{-1}\qquad(n\ge N_k),
\]
and set \(\varrho_n=k^{-1}\) when \(N_k\le n<N_{k+1}\).  Then
\(\varrho_n\downarrow0\) and
\(\chi_n=o_{\Pbb}(\varrho_n)\).  Equations
\eqref{eq:finite-edge-scaled-residual}--
\eqref{eq:finite-edge-scaled-jacobian-variation}, the uniformly bounded
inverse of \(\mathcal J_n\), and Newton--Kantorovich give a unique zero in
the scaled ball of radius \(\varrho_n\).  Its correction has input norm
\(o_{\Pbb}(\varrho_n)\), and therefore
\[
 |\tau_n-\tau^0|=o_{\Pbb}(\bar\delta\sqrt L/n),
 \qquad
 \max_i|z_{n,i}|=o_{\Pbb}(r_n),
 \qquad
 \max_i|m_{n,i}-\widetilde m_i^0|=o_{\Pbb}(\bar\delta).
\]
Together with the definition of the base point, these estimates give
\eqref{eq:finite-edge-contact-asymptotics} and
\eqref{eq:finite-edge-mass-asymptotics}.  Notice in particular that the
central block is \(-I+o_{\Pbb}(1)\).  The Hessian
\(\nabla^2B_n(0)\approx-\eps_nI\) is the Hessian of the unfitted bump
field; it is not the Jacobian obtained when the central atom and all
fitted denominators move together.

Stationarity in \(\tau\) and \(z_i\) gives the gradient contact equations,
while stationarity in \(m_i\) makes all displayed certificate values
equal.  Since
\(\int\Gamma_{\widetilde\Pi}\,d\widetilde\Pi=1\), their common value is
one.  For \(j\notin A\),
\[
 M_j(\widetilde\Pi)
 \ge
 \left(1-\frac{\sum_i m_{n,i}}n\right)k_{\tau_n}(j),
\]
so \(\max_{j\notin A}(\log c_j)_+=O_{\Pbb}(\bar\delta L/n)\).
The own fitted value and
\eqref{eq:finite-edge-mass-asymptotics} give
\(\log c_i=-\delta v_i+o_{\Pbb}(\bar\delta)\).

Finally, separate the \(i\)th observation in the certificate:
\[
 \Gamma_{\widetilde\Pi}(q_i+y)
 =
 c_i e^{\delta v_i}e^{-|y|^2/2}
 +\frac1n\sum_{j\ne i}c_j
   e^{(q_i+y)\cdot q_j-|q_i+y|^2/2}.
\]
The positive fitted-weight bound and the absolute envelopes imply
\eqref{eq:finite-edge-certificate-c2}.  The own Gaussian has Hessian
\(-I\) at its peak, so the certificate is strictly log-concave on a
smaller patch containing its stationary contact.  For
\(i\in\mathcal U\setminus A\), its adjusted own height is at most
\(\exp\{-\eta\delta+O_{\Pbb}(\bar\delta L/n)\}\), and the punctured
background is \(o_{\Pbb}(\bar\delta)\).  This proves the inactive-patch
claim and completes the proof.
\end{proof}

\subsubsection{Support identity when \(d_n\log n=o(n^2)\)}

The finite-contact lemma is deterministic. We now verify its hypotheses
throughout the range \(d_nL_n=o(n^2)\) and express the resulting support
identity in the exact coordinate of Theorem~\ref{thm:main-support-law}.

\begin{proposition}
\label{prop:buffered-support-below-simplex}
Suppose
\[
 d_n\longrightarrow\infty,
 \qquad d_nL_n=o(n^2),
 \qquad s_n^\star=O(1).
\]
For every fixed \(C,J<\infty\) and \(\eta>0\),
\begin{equation}
\label{eq:buffered-support-below-simplex}
 \Pbb\left\{
  \mathcal B_n(C,J,\eta),\quad
  \widehat K_n\ne1+\#\{i:\xi_{n,i}^\star>0\}
 \right\}\longrightarrow0.
\end{equation}
On \(\mathcal B_n(C,J,\eta)\), the NPMLE is also unique apart from an event
whose probability tends to zero.
\end{proposition}

\begin{proof}
Lemma~\ref{lem:exact-old-edge-coordinate} allows us to use
\(v_{n,i}=\xi_{n,i}-s_n\) in place of \(\xi_{n,i}^\star\).  On a slightly
smaller buffer the two active sets agree.  Put
\(A_n=\{i:v_{n,i}>0\}\); throughout the proof \(|A_n|\le J\).

We first describe the finite perturbation common to all the dimension
ranges.  The centered directions are not conditionally independent, so we
embed \(A_n\) in an uncentered buffered cloud.  The uniform
centered--uncentered edge-coordinate estimates in
Lemmas~\ref{lem:centered-radial-max} and
\ref{lem:centered-radial-max-high-d} show that, with probability tending
to one, every \(i\in A_n\) belongs to
\[
 \mathcal V_n=\{i:\xi_{n,i}^\circ-s_n>-\eta/2\}.
\]
This cloud has \(O_{\Pbb}(1)\) points. Conditional on its uncentered
radii, the directions are independent and uniform. Since \(d_n\to\infty\),
spherical concentration gives the following separation for the
uncentered score vectors. The common empirical-centering shift is
\(o_{\Pbb}(\sqrt{L_n})\), so it does not change their leading pairwise inner
products or distances. Consequently,
\begin{equation}
\label{eq:finite-edge-angular-separation}
 \max_{i\ne j\in A_n}
 \frac{|q_i\cdot q_j|}{L_n}=o(1),
 \qquad
 |q_i-q_j|^2=4L_n\{1+o(1)\}.
\end{equation}
The pairwise estimates in \eqref{eq:finite-edge-angular-separation} also hold
between an active point and any fixed number of other points in the upper
radial cloud.  Hence a neighborhood of \(q_i\)
sees its own bump at order one and every other retained upper bump at
\(n^{-1+o(1)}\).

Let \(\delta_n=a_n/\gamma_n\) and
\(\bar\delta_n=1\wedge\delta_n\).  The regime expansions in
Section~\ref{sec:regime-map} give \(\delta_n\le2\) for all large \(n\).
They also verify the scale assumptions in
Lemma~\ref{lem:uniform-finite-edge-expansion}. If
\(d_n=O(L_n)\), then \(\bar\delta_n\) is bounded away from zero.  If
\(L_n=o(d_n)\), then
\[
 \delta_n=\sqrt{L_n/d_n}\{1+o(1)\},
 \qquad
 n\bar\delta_n
 =\frac{L_n}{\sqrt{d_nL_n/n^2}}\{1+o(1)\}\longrightarrow\infty,
\]
because \(d_nL_n=o(n^2)\).  Thus
\(n\bar\delta_n\to\infty\) in the slowly growing, direct-truncation,
crossover, and entropy ranges alike.
After replacing \(\eta\) by a smaller fixed constant, the coordinate
comparison and the buffer imply that every nonactive index has
\(v_{n,i}\le-\eta\).  We may therefore take \(\mathcal U=A_n\) in
Lemma~\ref{lem:uniform-finite-edge-expansion}.
We verify the hypotheses of
Lemma~\ref{lem:uniform-finite-edge-expansion} uniformly over the finitely
many active indices.

Fix a small constant \(r_0<1/4\).  In the slowly growing, direct
truncation, and crossover ranges, let \(R_n=r_0\).  In the entropy range,
let \(R_n=L_n^{-1/2}\).  Define the available-radius envelopes
\[
 \widetilde e_{k,n}
 =
 \max_{i\in A_n}\mathcal E_{i,k}(2R_n),
 \qquad 0\le k\le3,
\]
and put
\[
 \Lambda_n^{\rm loc}
 =\widetilde e_{1,n}\vee\frac{\bar\delta_nL_n}{n},
 \qquad
 U_n=R_n\wedge\sqrt{\bar\delta_n}.
\]
The earlier field estimates imply, in each range, that
\begin{equation}
\label{eq:newton-radius-separation}
 \frac{\Lambda_n^{\rm loc}}{U_n}=o_{\Pbb}(1),
 \qquad
 \widetilde e_{0,n}=o_{\Pbb}(\bar\delta_n),
 \qquad
 \widetilde e_{2,n}+\widetilde e_{3,n}=o_{\Pbb}(1).
\end{equation}

In the slowly growing range, the proof of
Lemma~\ref{lem:slow-punctured-c3} bounds the sum of the absolute
derivatives of the lower cloud.  The remaining upper cloud has at most
\(e^{2\sqrt{d_n}}\) points, and angular separation bounds its first three
derivative envelopes by
\[
 e^{2\sqrt{d_n}}L_n^{3/2}e^{-2L_n+o(L_n)}=o(1).
\]
Thus \(\widetilde e_{0,n}=o_{\Pbb}(1)\) and
\(\sum_{k=1}^3\widetilde e_{k,n}=o_{\Pbb}(1)\).  Since
\(\bar\delta_n=1+o(1)\), \(U_n\) is bounded away from zero, and
\(L_n/n=o(1)\), these estimates give \eqref{eq:newton-radius-separation}.

In the direct truncation range, positivity of the Gaussian bumps and the
proof of Proposition~\ref{prop:moderate-low-cloud} give
\begin{equation}
\label{eq:direct-available-envelope}
 \widetilde e_{k,n}
 \le
 O_{\Pbb}(\delta_nL_n^{-10+k/2})+e^{-L_n+o(L_n)},
 \qquad 0\le k\le3.
\end{equation}
When \(\bar\delta_n\) vanishes, the regime condition
\(d_n(\log d_n)^2=o(L_n^3)\) and the superlogarithmic expansion imply
\(\bar\delta_n\ge L_n^{-1+o(1)}\).  Hence
\(\widetilde e_{1,n}/U_n=o_{\Pbb}(1)\) and
\((\bar\delta_nL_n/n)/U_n=o(1)\). The bounds in
\eqref{eq:direct-available-envelope} for \(k=0,2,3\) prove the other two assertions in
\eqref{eq:newton-radius-separation}.

In the crossover range, \eqref{eq:crossover-separation} gives
\begin{equation}
\label{eq:crossover-absolute-envelope}
 \widetilde e_{k,n}
 \le
 L_n^{k/2}e^{-L_n/2+O(\sqrt{L_n})}
 =o(\delta_n^m)
 \qquad(0\le k\le3)
\end{equation}
for every fixed \(m\).  Since
\(\bar\delta_n\asymp(\log L_n)/L_n\) and
\(U_n\asymp\sqrt{\bar\delta_n}\), the superpolynomial bound and
\(\sqrt{L_n\log L_n}/n=o(1)\) give
\eqref{eq:newton-radius-separation}.

Finally, in the entropy range put
\[
 \mu_n^{\rm off}=\max_{i\ne j}|q_i\cdot q_j|.
\]
The pairwise Gram bound gives
\begin{equation}
\label{eq:entropy-local-absolute-envelope}
 \widetilde e_{k,n}
 \le
 \frac{e^{\mu_n^{\rm off}+O_{\Pbb}(1)}L_n^{k/2}}n,
 \qquad 0\le k\le3.
\end{equation}
Writing \(A_n^{\rm off}=e^{\mu_n^{\rm off}}/n\), the rates in
Lemma~\ref{lem:random-off-diagonal} imply
\[
 A_n^{\rm off}=o_{\Pbb}(\delta_n),\qquad
 A_n^{\rm off}L_n=o_{\Pbb}(1),\qquad
 A_n^{\rm off}L_n^{3/2}=o_{\Pbb}(1).
\]
Since \(U_n^{-1}\le\sqrt{L_n}+\bar\delta_n^{-1/2}\),
\(A_n^{\rm off}\sqrt{L_n/\bar\delta_n}
=\{(A_n^{\rm off}/\bar\delta_n)(A_n^{\rm off}L_n)\}^{1/2}
=o_{\Pbb}(1)\).  It follows that
\(\widetilde e_{1,n}/U_n=o_{\Pbb}(1)\); also
\((\bar\delta_nL_n/n)/U_n=o(1)\), since it is at most
\(\bar\delta_nL_n^{3/2}/n+\sqrt{\bar\delta_n}L_n/n\).  The three rates
for \(A_n^{\rm off}\), applied to
\eqref{eq:entropy-local-absolute-envelope}, prove the remaining assertions
in \eqref{eq:newton-radius-separation}.

We now choose the Newton scale.  From
\eqref{eq:newton-radius-separation}, choose integers \(M_k\uparrow\infty\)
so that
\[
 \Pbb\left\{\Lambda_n^{\rm loc}/U_n>k^{-2}\right\}\le k^{-1}
 \qquad(n\ge M_k),
\]
and set \(c_n=k^{-1}\) for \(M_k\le n<M_{k+1}\).  Then \(c_n\downarrow0\)
and
\[
 \Lambda_n^{\rm loc}=o_{\Pbb}(c_nU_n).
\]
Set \(r_n=c_nU_n\).  Then
\[
 r_n\le R_n,\qquad
 \widetilde e_{1,n}=o_{\Pbb}(r_n),\qquad
 \frac{\bar\delta_nL_n}{n}=o(r_n),\qquad
 \frac{r_n^2}{\bar\delta_n}\le c_n^2\longrightarrow0.
\]
Because the envelopes decrease with the patch radius, the three inequalities
in \eqref{eq:newton-radius-separation} verify
\eqref{eq:finite-edge-envelope-assumptions}.  Moreover,
\[
 r_n\widetilde e_{1,n}=o_{\Pbb}(r_n^2)=o_{\Pbb}(\bar\delta_n),
 \qquad
 r_n^2\widetilde e_{2,n}+r_n^3\widetilde e_{3,n}
 =o_{\Pbb}(\bar\delta_n),
\]
so every condition in
\eqref{eq:finite-edge-higher-envelope-assumptions} holds.

Lemma~\ref{lem:uniform-finite-edge-expansion} therefore constructs a
measure \(\widetilde\Pi\) with one central contact \(\tau_n\) and one
contact \(w_{n,i}\) in each active patch.  Its edge weights satisfy
\begin{equation}
\label{eq:below-simplex-contact-asymptotics}
 |\tau_n|=O_{\Pbb}(\bar\delta_n\sqrt{L_n}/n),\qquad
 w_{n,i}=q_i+o_{\Pbb}(r_n),\qquad
 n\pi_{n,i}
 =1-e^{-\delta_nv_{n,i}}+o_{\Pbb}(\bar\delta_n).
\end{equation}
The buffer keeps every displayed mass positive.  In the entropy range,
retaining the explicit right side of
\eqref{eq:entropy-local-absolute-envelope} in the contact-map proof gives
\begin{equation}
\label{eq:entropy-local-contact-rate}
 n\pi_{n,i}
 =
 1-e^{-\delta_nv_{n,i}}
 +O_{\Pbb}\left\{
   \frac{e^{\mu_n^{\rm off}}}n+\frac{\delta_nL_n}{n}
 \right\}.
\end{equation}

It remains to check that the finite perturbation has not created a contact
elsewhere.  If \(c_j\) denotes the reciprocal fitted value at observation
\(j\), then
\begin{equation}
\label{eq:below-simplex-weight-perturbation}
 (\log c_j)_+
 =
 O_{\Pbb}\left\{\frac{\bar\delta_nL_n}{n}\right\}
 \quad(j\notin A_n),
 \qquad
 \log c_i=-\delta_nv_{n,i}+o_{\Pbb}(\bar\delta_n)
 \le-c_\eta\bar\delta_n
 \quad(i\in A_n).
\end{equation}
In the entropy range the contact equation gives the sharper cancellation
\begin{equation}
\label{eq:entropy-contact-cancellation}
 \max_{i\in A_n}
 \left|\delta_nv_{n,i}+\log c_i\right|
 +\max_{j\notin A_n}|\log c_j|
 =
 O_{\Pbb}\left\{
  \frac{e^{\mu_n^{\rm off}}}n+\frac{\delta_nL_n}{n}
 \right\}
 =n^{-1+o_{\Pbb}(1)}.
\end{equation}
We first control the inactive denominators without using the contact
equation.  If \(m=\sum_{i\in A_n}n\pi_{n,i}=O_{\Pbb}(\delta_n)\), then
\[
 M_j
 =
 \left(1-\frac mn\right)
 e^{\tau_n\cdot q_j-|\tau_n|^2/2}
 +\frac1n\sum_{i\in A_n}n\pi_{n,i}
 e^{w_{n,i}\cdot q_j-|w_{n,i}|^2/2}.
\]
Here
\(|\tau_n|=O_{\Pbb}(\delta_n\sqrt{L_n}/n)\), so the central term differs
from one by \(O_{\Pbb}(\delta_nL_n/n)\).  For \(j\notin A_n\), each edge
response in the last sum is at most
\(n^{-1}e^{\mu_n^{\rm off}+o_{\Pbb}(1)}\).  Consequently,
\[
 \max_{j\notin A_n}|M_j-1|
 =
 O_{\Pbb}\left\{
  \frac{\delta_nL_n}{n}
  +\frac{\delta_ne^{\mu_n^{\rm off}}}{n^2}
 \right\},
\]
which proves the inactive part of
\eqref{eq:entropy-contact-cancellation}.

We may now evaluate the exact level-one contact equation in the \(i\)th
edge patch.  If \(w_i\) is its contact and \(E_i\) is the contribution of
all observations other than \(i\) to the certificate there, then the own
term gives the exact identity
\begin{equation}
\label{eq:entropy-exact-contact-cancellation}
 \log M_i
 =
 \delta_nv_{n,i}-\frac{|w_i-q_i|^2}{2}-\log(1-E_i).
\end{equation}
The derivative envelope
\eqref{eq:entropy-local-absolute-envelope}, the \(C^2\) certificate
expansion in Lemma~\ref{lem:uniform-finite-edge-expansion}, and the local
stationarity equation give
\[
 |w_i-q_i|^2
 =O_{\Pbb}\left\{
  \frac{e^{2\mu_n^{\rm off}}L_n}{n^2}
  +\frac{\delta_n^2L_n^2}{n^2}
 \right\},
 \qquad
 E_i
 =O_{\Pbb}\left\{
  \frac{e^{\mu_n^{\rm off}}}n+\frac{\delta_nL_n}{n}
 \right\}.
\]
Indeed, a retained upper response is at most
\(e^{\mu_n^{\rm off}}/n\), and the
central displacement and lower fitted values contribute
\(O_{\Pbb}(\delta_nL_n/n)\).  Substitution in
\eqref{eq:entropy-exact-contact-cancellation} proves the active part of
\eqref{eq:entropy-contact-cancellation}.  These estimates are uniform
because \(|A_n|\le J\), and the argument is noncircular: the inactive
denominator bound was obtained before the certificate contribution
\(E_i\) was estimated.
The positive part in
\eqref{eq:below-simplex-weight-perturbation} is \(o_{\Pbb}\) of the
smallest gap used in each one-atom exclusion proof.  The following
four-region argument verifies these comparisons while proving the
weighted exclusion rather than applying the unweighted results formally.
Write
\[
 \Gamma_{\widetilde\Pi}(w)
 =
 \Gamma_I(w)+\sum_{i\in A_n}G_i(w),
 \qquad
 G_i(w)=\frac{c_i}{n}e^{w\cdot q_i-|w|^2/2},
\]
where \(I=A_n^c\) and
\[
 \Gamma_I(w)
 =\frac1n\sum_{j\in I}c_j
   e^{w\cdot q_j-|w|^2/2},
 \qquad
 D_I(w)
 =\frac1n\sum_{j\in I}
   e^{w\cdot q_j-|w|^2/2}.
\]
If
\[
 \omega_n
 =\max_{j\in I}(\log c_j)_+,
\]
then, pointwise and after every truncation used in the cited proofs,
\begin{equation}
\label{eq:weighted-inactive-domination}
 \Gamma_I(w)\le e^{\omega_n}D_I(w).
\end{equation}
Thus every Bernstein or net upper bound for the inactive field changes by
at most \(O(\omega_n)\).

There are four regions.  First consider a fixed small ball around an active
contact.  The own term satisfies the exact relation
\[
 G_i(w_{n,i})=1-E_i,
\]
where \(E_i\) is the remaining certificate contribution, rather than
being one by itself.  The \(C^2\) representation
\eqref{eq:finite-edge-certificate-c2} gives
\(\nabla^2\log\Gamma_{\widetilde\Pi}\preceq-cI\) throughout this ball.
Since the contact is stationary and has value one, it is the unique
level-one point there.

Next take a corresponding ball around an inactive upper point \(q_j\).
Its adjusted own peak is at most
\[
 \frac{c_j}{n}e^{|q_j|^2/2}
 \le
 \exp\{\omega_n+\delta_nv_{n,j}\}
 \le1-c_\eta\bar\delta_n.
\]
The punctured \(C^3\) estimates
\eqref{eq:moderate-low-cloud},
\eqref{eq:entropy-local-absolute-envelope},
\eqref{eq:slow-punctured-c3}, and
\eqref{eq:crossover-absolute-envelope},
in their respective ranges, make all remaining terms
\(o_{\Pbb}(\bar\delta_n)\).  Hence an inactive upper patch contains no
contact.

In the central region, the weighted Hessian differs from the unweighted
one by
\[
 O_{\Pbb}(\omega_nL_n)
 +O_{\Pbb}\left\{\frac{\bar\delta_nL_n^2}{n}\right\}
 =o_{\Pbb}(\eps_n).
\]
To obtain this bound, write each fitted value as \(M_j=1+x_j\).
The negative part of \(x_j\) is uniformly
\(O_{\Pbb}(\bar\delta_nL_n/n)\), while
\[
 \frac1n\sum_j(x_j)_+
 \le
 C\frac{\bar\delta_n}{n}
 +O_{\Pbb}\left\{\frac{\bar\delta_nL_n}{n}\right\};
\]
the first term follows by averaging each edge response and using
\(B_n(w_{n,i})=O_{\Pbb}(1)\). Multiplication by
\(\max_j|q_j|^2=O_{\Pbb}(L_n)\) gives the displayed covariance
perturbation.
It is therefore negative definite in the microscopic ball around the
translated central contact.  On the surrounding annuli,
\eqref{eq:weighted-inactive-domination} raises the field by at most
\(O_{\Pbb}(\omega_n)\).  We compare the active profiles with the annular
gap explicitly.  Uniformly for \(i\in A_n\),
\[
 G_i(w)
 \le
 \exp\left\{-\frac12(|q_i|-|w|)^2+o_{\Pbb}(1)\right\},
 \qquad
 |q_i|=\sqrt{2L_n}\{1+o_{\Pbb}(1)\}.
\]
Let \(g_n(r)=1-e^{-\eps_nr^2/2}\), and let \(g_{0,n}\) be the smallest
gap at the inner endpoint of the annular argument.  If
\(\eps_nr^2\le1\), then
\[
 g_n(r)\ge c\eps_nr^2\ge c g_{0,n},
 \qquad
 G_i(w)\le n^{-1+o(1)}=o(g_{0,n}).
\]
The last comparison is \(n^{-1}=o(\eps_n^3/L_n^3)\) in the slow and
direct truncation ranges and is immediate in the crossover range. If
\(\eps_nr^2\ge1\) and \(r\le|q_i|/2\), then \(g_n(r)\ge c\) and
\(G_i(w)\le n^{-1/4+o(1)}\).  Hence
\begin{equation}
\label{eq:active-central-corridor}
 \sum_{i\in A_n}G_i(w)=o_{\Pbb}\{g_n(|w|)\}
 \quad
 (r_{0,n}\le|w|\le\tfrac12\sqrt{2L_n}).
\end{equation}
Together with the annular estimates established separately in the slow,
direct-truncation, crossover, and entropy regimes, this proves that the
central contact is the only contact in the inner region.

Finally consider the complement of the central and upper patches.
Choose their fixed radius \(R=R(J)\) so large that
\(Je^{-R^2/2}<1/20\); the patches are disjoint by angular separation.
On the annulus between the implicit-function neighborhood and radius
\(R\), the punctured \(C^2\) bound and the exact level-one contact show
that the Gaussian profile generated by the nearest active observation has
its unique maximum at the constructed contact.  In the slow range this
uniform statement is precisely Lemma~\ref{lem:slow-punctured-c3}; in the
direct-truncation and crossover ranges it follows from the local \(C^2\)
estimates used to construct the active contacts.
The inactive part is again controlled by
\eqref{eq:weighted-inactive-domination}.  Conditional angular separation
implies that at most one upper profile is nearest; outside its patch that
profile loses a fixed Gaussian factor, while the sum of all other upper
profiles is \(n^{-1+o_{\Pbb}(1)}\) in the direct truncation range,
\(e^{-L_n/2+o(L_n)}\) in the crossover range, and \(o_{\Pbb}(1)\) in the
slow range by Lemma~\ref{lem:slow-punctured-c3}.  Thus the outer-region
decompositions from the unweighted proofs remain valid after the fitted
denominators are inserted: equation
\eqref{eq:weighted-inactive-domination} controls the inactive field, and
the active tails retain their original bounds.  These estimates exclude
every remaining level-one point.
More concretely, in the crossover proof
\eqref{eq:crossover-outer-net} bounds the clipped inactive field by
\(0.9+o_{\Pbb}(1)\); choose \(R\) so that the active tails total less
than \(0.05\).  In the slow proof, choose \(R\) so that those tails are
smaller than \(\rho/4\), where \(\rho\) is the fixed slack in
Lemma~\ref{lem:slow-outside-shell}; the fitted-denominator error is also
\(o_{\Pbb}(\rho)\). The direct-truncation outer estimate has
an \(o_{\Pbb}(\delta_n)\) inactive background and is stronger still.

In the entropy range,
Lemma~\ref{lem:weighted-entropy-localization}, applied with
\(\beta_j=\log c_j\), directly localizes every level-one contact to the
central patch or an active upper patch;
\eqref{eq:entropy-contact-cancellation} verifies all of the lemma's weight
hypotheses.  The strict log-concavity asserted there leaves only the
constructed contacts.

The fitted-denominator error and the tails of the active profiles are
smaller than the KKT gaps in each unweighted localization argument.  More
precisely,
\[
 \frac{L_n}{n}=o\left(\frac{\eps_n^3}{L_n^3}\right)
 \quad\text{in the slow range},\qquad
 \frac{\delta_nL_n}{n}+\delta_n\mathfrak j_n=o(g_{0,n})
 \quad\text{in the direct range},
\]
where \(g_{0,n}\) is from
Proposition~\ref{prop:moderate-inner-field}; in the crossover range,
\(\delta_nL_n/n=o(L_n^{-3})\) and
\(e^{-L_n/2+o(L_n)}=o(\delta_n^m)\) for every fixed \(m\).
Thus the weighted estimates retain the gaps and negative curvature used in
all four regions.

This verifies the global KKT inequality for the constructed measure. Strict
concavity of the log likelihood makes its fitted vector unique. The contact likelihood
vectors are linearly independent: their active-coordinate block has
diagonal entries \(n^{1+o(1)}\) and off-diagonal entries \(n^{o(1)}\),
while the central vector is bounded.  Since
 \(n^{-1}\sum_jk_{w_{n,i}}(j)=B_n(w_{n,i})=O_{\Pbb}(1)\) for each of the
finitely many edge columns, their combined column sum is \(O_{\Pbb}(n)\).
Choose a deterministic \(K_n\to\infty\) so slowly that \(K_n=o(n)\).
Then
\[
 \#\left\{j:\max_{i\in A_n}k_{w_{n,i}}(j)>K_n\right\}
 \le
 \frac1{K_n}\sum_{i\in A_n}\sum_{j=1}^n k_{w_{n,i}}(j)
 =o_{\Pbb}(n).
\]
Since \(|A_n|\le J\), with probability tending to one there is an inactive
observation at which every edge entry is at most \(K_n=o(n)\).  Adjoin that
observation as one extra row; its central entry is
\(1+o_{\Pbb}(1)\).  After scaling the edge columns by \(n^{-1}\), their
active-coordinate block is asymptotically diagonal with diagonal bounded
away from zero, while all edge entries in the added row are \(o(1)\).
The resulting square minor is therefore nonsingular.  Every
maximizing mixing measure must therefore use all and only the
\(1+|A_n|\) contacts.

Finally, pass to a subsequence on which
\(\mathfrak r_n=d_n(\log d_n)^2/L_n^3\) tends to zero, stays bounded above
and away from zero, or tends to infinity.  These cases invoke,
respectively, the slow/direct estimates (split at
\(d_n=(\log L_n)^2\)), the crossover estimates, and the entropy estimates.
They exhaust \(d_nL_n=o(n^2)\) and prove
\eqref{eq:buffered-support-below-simplex}.
\end{proof}

\subsubsection{Ordered contacts in the simplex regime}

When \(d_nL_n/n^2\) is bounded below, the centered sample is close to a
regular simplex and the isolated-bump model no longer applies.  A radial
crossing then creates a contact near an ordered Potts state.  Write
\[
 \chi_n=\frac{n-2}{n-1},
 \qquad
 p_i^{(i)}=\frac{n-1}{n},
 \qquad
 p_j^{(i)}=\frac1{n(n-1)}\quad(j\ne i).
\]
The identities \(\sum_jq_j=0\) and
\eqref{eq:potts-radial-threshold} give, exactly,
\begin{equation}
\label{eq:ordered-potts-excess}
 \sum_jp_j^{(i)}q_j=\chi_nq_i,
 \qquad
 \frac12\left|\sum_jp_j^{(i)}q_j\right|^2
 -K(p^{(i)}\Vert \mathbf u_n)
 =\chi_n^2\delta_n^\star \xi_{n,i}^\star.
\end{equation}
Thus the sign of the ordered-state free energy is the sign of the exact
radial excess, with no centering approximation.

\begin{proposition}
\label{prop:buffered-simplex-support}
Let \(n\ge3\), let \(q_1,\ldots,q_n\) be a centered triangular array, and
write
\[
 R_i=\frac{|q_i|^2}{2},
 \qquad
 u_i=\frac{R_i-R_n^\star}{\delta_n},
 \qquad
 G_n=(q_i\cdot q_j)_{i,j\le n}.
\]
Suppose \(\delta_n\downarrow0\), \(\delta_n=O(L_n/n)\), and, for fixed
\(C,J<\infty\) and \(\eta>0\),
\[
 \max_i u_i\le C,
 \qquad
 \#\{i:u_i>-\eta\}\le J,
 \qquad
 \min_i|u_i|>\eta.
\]
Assume also that
\begin{equation}
\label{eq:buffered-simplex-assumptions}
 \max_{i\ne j}|q_i\cdot q_j|=o(L_n^{-1}),
 \qquad
 \|G_n\|_{\rm op}=O(L_n),
 \qquad
 \frac{C_n}{\delta_nL_n^2}\longrightarrow0,
 \qquad
 \frac{C_n^2}{n\delta_n}\longrightarrow0,
\end{equation}
where
\[
 C_n=\max_{i\ne j}
 \left|q_i\cdot q_j+\frac{|q_i|^2}{n-1}\right|.
\]
Then, for all sufficiently large \(n\), every NPMLE has
\begin{equation}
\label{eq:buffered-simplex-count}
 |\supp(\Pi)|=1+\#\{i:u_i>0\}.
\end{equation}
There is one central contact and, for each \(i\) with \(u_i>0\), one
contact \(w_i=\chi_nq_i+o(1)\); its weight satisfies
\begin{equation}
\label{eq:buffered-simplex-mass}
 n\pi_i=\frac{n}{n-1}\delta_nu_i+o(\delta_n).
\end{equation}
All remainders are uniform when \(C,J,\eta\) in the buffer are fixed.
Moreover, the NPMLE is unique.
\end{proposition}

\begin{proof}
Put \(A=\{i:u_i>0\}\), so \(|A|\le J\), and set
\(w_i^0=\chi_nq_i\) for \(i\in A\).  The buffer and the definition of
\(R_i\) imply
\begin{equation}
\label{eq:terminal-simplex-derived-rates}
 \max_i|q_i|=O(\sqrt{L_n}),
 \qquad
 C_n=o(L_n^{-1}),
 \qquad
 \delta_nL_n=o(1),
 \qquad
 \frac{C_n^2}{n}=o(\delta_n).
\end{equation}
Indeed, the first assertion follows from
\(R_i\le R_n^\star+C\delta_n=O(L_n)\), and the second follows from the
first condition in \eqref{eq:buffered-simplex-assumptions} and
\(|q_i|^2/(n-1)=O(L_n/n)\).  The last two assertions follow directly from
the assumed bounds on \(\delta_n\) and \(C_n\).

If \(A=\varnothing\), then
\(\max_iR_i\le R_n^\star-\eta\delta_n\).  Proposition
\ref{prop:local-simplex-criterion}, with \(g_n=\eta\delta_n\), applies by
\eqref{eq:buffered-simplex-assumptions}.  Its proof is strict away from
\(p=\mathbf u_n\), so the one-atom certificate equals one only at the
origin.  The KKT conditions then force every NPMLE to be \(\delta_0\), and
the conclusion follows.  We henceforth assume that \(A\ne\varnothing\).

We begin by locating the maxima of the one-atom certificate.  If
\[
 e_{ij}=q_i\cdot q_j+\frac{|q_i|^2}{n-1},
\]
then \(|e_{ij}|\le C_n\), and direct substitution gives
\begin{align}
\label{eq:potts-response-expansion}
 k_{w_i^0}(i)
 &= (n-1)\exp\{(2\chi_n-\chi_n^2)\delta_nu_i\},\notag\\
 k_{w_i^0}(j)
 &= \frac1{n-1}
    \exp\{-(\chi_n^2+2\chi_n/(n-1))\delta_nu_i
                 +\chi_ne_{ij}\},\qquad j\ne i.
\end{align}
At \(u_i=0\) and \(e_{ij}=0\), these responses are \(n-1\) and
\((n-1)^{-1}\), respectively, and their average is one.  Moreover,
\(\sum_{j\ne i}e_{ij}=0\), by centering.  The linear row error therefore
cancels when the off-diagonal responses are summed.  Expanding
\eqref{eq:potts-response-expansion}, using \(C_n=o(1)\), gives
\begin{equation}
\label{eq:potts-local-height}
 \log B_n(w_i^0)
 =\chi_n^2\delta_nu_i
  +O_{C,J}\left(\delta_n^2+\frac{C_n^2}{n}\right)
 =\chi_n^2\delta_nu_i+o(\delta_n).
\end{equation}
The coefficient can also be read directly from the exact dual identity
\eqref{eq:ordered-potts-excess}.

The Gibbs weights at \(w_i^0\) differ from \(p^{(i)}\) by
\(O_{C,J}\{(C_n+\delta_n)/n\}\) in total variation.  Differentiating the
log-sum-exp representation therefore gives, uniformly for \(i\in A\),
\begin{align}
\label{eq:potts-local-derivatives}
 |\nabla\log B_n(w_i^0)|
 &\le K_{C,J}\frac{\sqrt{L_n}}n(C_n+\delta_n),\notag\\
 \nabla^2\log B_n(w_i^0)
 &=-I+O(L_n/n)+o(1)=-I+o(1).
\end{align}
The same estimate holds uniformly on \(B(w_i^0,L_n^{-2})\): moving by
\(L_n^{-2}\) changes every exponent by \(O(L_n^{-3/2})\), while the total
Gibbs mass outside coordinate \(i\) remains \(O(n^{-1})\).  The
Newton--Kantorovich theorem now gives a unique critical point \(\bar w_i\)
in that ball, with
\begin{equation}
\label{eq:quantitative-potts-location}
 |\bar w_i-w_i^0|
 \le K_{C,J}\frac{\sqrt{L_n}}n(C_n+\delta_n).
\end{equation}
It is a strict local maximum, and Taylor's theorem gives
\begin{equation}
\label{eq:potts-corrected-height}
 h_i:=\log B_n(\bar w_i)
 =\chi_n^2\delta_nu_i+\varepsilon_{n,i},
 \qquad
 \max_{i\in A}|\varepsilon_{n,i}|=o(\delta_n).
\end{equation}
More explicitly, the remainder is bounded by a constant times
\[
 \delta_n^2+\frac{C_n^2}{n}
 +\frac{L_n(C_n+\delta_n)^2}{n^2},
\]
which is \(o(\delta_n)\) by
\eqref{eq:terminal-simplex-derived-rates} and
\eqref{eq:buffered-simplex-assumptions}.  In particular,
\(h_i\ge c_\eta\delta_n\) for every \(i\in A\) and all sufficiently large
\(n\).

For \(i\in A\), let
\(v_i=(k_{\bar w_i}(1),\ldots,k_{\bar w_i}(n))\).  Equations
\eqref{eq:potts-response-expansion}--\eqref{eq:potts-local-derivatives}
give
\begin{equation}
\label{eq:potts-response-gram}
 A_{ij}^{(0)}
 :=\frac1{n^2}\sum_{r=1}^n\{v_i(r)-1\}\{v_j(r)-1\}
 =\beta_n^{\rm P}\mathbf1_{\{i=j\}}+e_{n,ij},
 \qquad
 \max_{i,j\in A}|e_{n,ij}|=o(1),
\end{equation}
where
\[
 \beta_n^{\rm P}=\frac{(n-2)^2}{n(n-1)}.
\]
For instance, at the unperturbed ordered state the diagonal expression is
exactly
\[
 \frac1{n^2}\left\{(n-2)^2
 +(n-1)\left(\frac1{n-1}-1\right)^2\right\}=\beta_n^{\rm P},
\]
whereas the cross expressions are \(O(n^{-1})\).  The perturbations in
\eqref{eq:potts-response-expansion} and
\eqref{eq:quantitative-potts-location} contribute \(o(1)\), uniformly over
the at most \(J\) active indices.

We next construct exact contacts and masses.  For
\(m=(m_i)_{i\in A}\) and
\(x=(\tau,(w_i)_{i\in A})\), define
\[
 \Pi_{x,m}
 =\left(1-\frac1n\sum_{i\in A}m_i\right)\delta_\tau
  +\frac1n\sum_{i\in A}m_i\delta_{w_i},
\]
and write \(\Gamma_{x,m}\) for its certificate.  Consider the stationarity
map
\[
 \mathcal S(x,m)
 =\left(
   \nabla\log\Gamma_{x,m}(\tau),
   \bigl(\nabla\log\Gamma_{x,m}(w_i)\bigr)_{i\in A}
  \right).
\]
At \(m=0\) and
\(x=x^0:=(0,(\bar w_i)_{i\in A})\), this map vanishes.  Equip the product
of location spaces with the maximum of the Euclidean block norms.  For
each fixed \(M<\infty\), the response table gives, uniformly on
\[
 |m|_\infty\le M\delta_n,
 \qquad |\tau|\le L_n^{-2},
 \qquad |w_i-\bar w_i|\le L_n^{-2},
\]
the bounds
\begin{equation}
\label{eq:terminal-location-map-bounds}
 \|D_x\mathcal S(x^0,0)^{-1}\|\le2,
 \quad
 \|D_x\mathcal S(x,m)-D_x\mathcal S(x^0,0)\|=o(1),
 \quad
 \|D_m\mathcal S(x,m)\|
 \le K_{C,J,M}\frac{\sqrt{L_n}}n.
\end{equation}
To verify the first two bounds, the diagonal location blocks are
\(-I+o(1)\) by \eqref{eq:potts-local-derivatives}, while the central block
is
\[
 \frac1n\sum_{r=1}^nq_rq_r^\top-I=-I+O(L_n/n).
\]
At zero mass, the derivative of an active equation with respect to the
central location is also negligible: if the central atom is at \(\tau\),
then
\[
 \log\Gamma_{\delta_\tau}(w)
 =\log B_n(w-\tau)-(w-\tau)\cdot\tau,
\]
so at \((\tau,w)=(0,\bar w_i)\),
\[
 D_\tau\nabla_w\log\Gamma_{\delta_\tau}(w)
 =-\nabla^2\log B_n(\bar w_i)-I=o(1).
\]
All other cross-location blocks vanish at zero mass and are
\(O_{C,J}(\delta_nL_n)=o(1)\) on the displayed neighborhood.  Finally,
one differentiation with respect to \(m_j\) produces a normalized product
of a central or ordered response and its first moment.  The response table
bounds this by \(K\sqrt{L_n}/n\), proving the last estimate in
\eqref{eq:terminal-location-map-bounds}.

The quantitative implicit-function theorem now supplies unique smooth
locations satisfying the stationarity equations.  Write
\[
 x(m)=\bigl(\tau(m),(w_i(m))_{i\in A}\bigr).
\]
They obey
\begin{equation}
\label{eq:terminal-location-map-rate}
 |\tau(m)|+\max_{i\in A}|w_i(m)-\bar w_i|
 \le K_{C,J,M}\frac{\sqrt{L_n}}n|m|_1.
\end{equation}
Define the logarithmic contact-height differences
\[
 \mathcal H_i(m)
 =\log\Gamma_{x(m),m}\{w_i(m)\}
  -\log\Gamma_{x(m),m}\{\tau(m)\}.
\]
Stationarity permits the envelope theorem.  At zero mass,
\(\mathcal H_i(0)=h_i\), and direct differentiation gives
\begin{equation}
\label{eq:terminal-mass-contact-jacobian}
 -\partial_{m_j}\mathcal H_i(0)
 =A_{ij}^{(0)}+o(1)
 =\beta_n^{\rm P}\mathbf1_{\{i=j\}}+o(1).
\end{equation}
The additional \(o(1)\), beyond
\eqref{eq:potts-response-gram}, includes the movement of the central
contact; by \eqref{eq:terminal-location-map-bounds} its contribution is
\(O(L_n/n)\).

The response averages in \eqref{eq:potts-response-gram} also control the
second derivative.  In particular,
for any \(i,j,k\in A\),
\[
 \frac1{n^3}\sum_{r=1}^n
  \{1+v_i(r)\}\{1+v_j(r)\}\{1+v_k(r)\}=O_{C,J}(1).
\]
The fitted denominators stay between \(1/2\) and \(2\) on
\(|m|_\infty\le M\delta_n\), so
\begin{equation}
\label{eq:terminal-mass-second-derivative}
 \sup_{|m|_\infty\le M\delta_n}
 \|D_m^2\mathcal H(m)\|=O_{C,J,M}(1).
\end{equation}
Consequently,
\begin{equation}
\label{eq:potts-local-program}
 \mathcal H(m)=h-A_nm+R_n(m),
 \qquad
 A_n=\beta_n^{\rm P}I_{|A|}+o(1),
 \qquad
 |R_n(m)|_\infty\le K_{C,J}|m|_\infty^2.
\end{equation}
Choose \(M\) larger than \(2C\).  Since \(u_i\in[\eta,C]\) on \(A\),
the map
\(m\mapsto A_n^{-1}\{h+R_n(m)\}\) is a contraction on a ball of radius
\(o(\delta_n)\) about \(A_n^{-1}h\).  It therefore has a unique fixed
point in \(|m|_\infty\le M\delta_n\), and
\begin{equation}
\label{eq:terminal-quantitative-mass}
 m_i
 =\frac{\chi_n^2}{\beta_n^{\rm P}}\delta_nu_i+o(\delta_n)
 =\frac{n}{n-1}\delta_nu_i+o(\delta_n).
\end{equation}
Every component is positive for all sufficiently large \(n\).  The
resulting measure \(\widetilde\Pi=\Pi_{x(m),m}\) has one stationary central
contact and one stationary contact in each active patch, all at the same
certificate level.  Since
\(\int\Gamma_{\widetilde\Pi}\,d\widetilde\Pi=1\), that level equals one.

We finally prove the global KKT inequality.  Put
\(c_r=M_r(\widetilde\Pi)^{-1}\).  The response expansion and the mass
formula imply
\begin{equation}
\label{eq:potts-fitted-weight-profile}
 |\log c_r|=O_{C,J}(\delta_nL_n/n)\quad(r\notin A),
 \qquad
 \log c_i=-\chi_n\delta_nu_i+o(\delta_n)\quad(i\in A).
\end{equation}
Indeed, \eqref{eq:terminal-location-map-rate} gives
\(|\tau(m)|=O_{C,J}(\delta_n\sqrt{L_n}/n)\) and
\(\max_i|w_i(m)-\bar w_i|=O_{C,J}(\delta_n\sqrt{L_n}/n)\).
For \(r\notin A\), the central response is
\[
 e^{\tau(m)\cdot q_r-|\tau(m)|^2/2}
 =1+O_{C,J}(\delta_nL_n/n),
\]
the total mass removed from the central atom is \(O_{C,J}(\delta_n/n)\),
and all active responses together contribute \(O_{C,J}(\delta_n/n^2)\).
Thus \(M_r=1+O_{C,J}(\delta_nL_n/n)\).  At an active observation \(i\),
\eqref{eq:potts-response-expansion} and
\eqref{eq:terminal-quantitative-mass} instead give
\[
 M_i
 =1+\frac{n-2}{n}\,
      \frac{n}{n-1}\delta_nu_i+o(\delta_n)
 =1+\chi_n\delta_nu_i+o(\delta_n).
\]
Taking negative logarithms proves
\eqref{eq:potts-fitted-weight-profile}.

For clarity, put
\[
 F_{\log c}(p)
 =\frac12\left|\sum_rp_rq_r\right|^2
  -K(p\Vert\mathbf u_n)+\sum_rp_r\log c_r.
\]
The weighted entropy dual \eqref{eq:weighted-entropy-dual} shows that
\(\Gamma_{\widetilde\Pi}(w)\ge1\) only if its Gibbs vector \(p=p(w)\)
satisfies \(F_{\log c}(p)\ge0\).  If
\(\max_rp_r\le1-L_n^{-2}\), the stability inequality
\eqref{eq:entropy-stability} may be applied after retaining the active
diagonal terms.  Indeed, writing \(r=p-\mathbf u_n\), their contribution is at
most
\[
 \sum_{i\in A}\delta_nu_i r_i^2
 \le\chi_n\sum_{i\in A}\delta_nu_i p_i+O_{C,J}(\delta_n/n^2),
\]
where the scalar inequality
\((p_i-1/n)^2\le\chi_np_i+n^{-2}\) holds for \(0\le p_i\le1\).
This is canceled by \eqref{eq:potts-fitted-weight-profile}.  The remaining
active-weight error is \(o(\delta_n)\), and the positive weight on
inactive coordinates is \(O(\delta_nL_n/n)=o(\delta_n)\).  Consequently,
\eqref{eq:entropy-stability} and
\eqref{eq:buffered-simplex-assumptions} give, for a deterministic sequence
\(\varpi_n\downarrow0\),
\[
 F_{\log c}(p)
 \le
 -\frac{c}{L_n}|p-\mathbf u_n|_1^2+\varpi_n\delta_n,
 \qquad \varpi_n\longrightarrow0.
\]
Thus nonnegative weighted free energy forces
\[
 |p-\mathbf u_n|_1=o\{\sqrt{\delta_nL_n}\}
 =o(L_n/\sqrt n).
\]
Because \(\sum_rq_r=0\), the corresponding
\(\mu(p)=\sum_rp_rq_r\) satisfies
\[
 |\mu(p)|
 \le \max_r|q_r|\,|p-\mathbf u_n|_1
 =o(L_n\sqrt{\delta_n})
 =o(L_n^{3/2}/\sqrt n)=o(1).
\]
For the Gibbs vector of an actual certificate point, the exact square
identity is
\begin{equation}
\label{eq:weighted-certificate-square-identity}
 \log\Gamma_{\widetilde\Pi}(w)
 =F_{\log c}\{p(w)\}-\frac12|w-\mu\{p(w)\}|^2.
\end{equation}
Therefore \(\Gamma_{\widetilde\Pi}(w)\ge1\) implies \(w=o(1)\) in the
diffuse case.

Suppose instead that \(p_i=1-y>1-L_n^{-2}\), and decompose \(p\) as in
the proof of Proposition~\ref{prop:local-simplex-criterion}: write
\(p_j=yr_j\) for \(j\ne i\) and
\[
 v_r=\sum_{j\ne i}\left(r_j-\frac1{n-1}\right)q_j,
 \qquad
 \mathbf u_{n-1}=\left(\frac1{n-1},\ldots,\frac1{n-1}\right).
\]
Applying the local-simplex decomposition from that proposition with
\(R_i=R_n^\star+\delta_nu_i\) gives
\begin{align}
\label{eq:weighted-potts-localization}
 &\frac12\left|\sum_rp_rq_r\right|^2-K(p\Vert \mathbf u_n)\notag\\
 &\qquad\le
 \delta_nu_i c_y^2
 -\{k_n(y)-R_n^\star c_y^2\}
 +y|c_y|C_n,
 \qquad c_y=1-\frac{ny}{n-1}.
\end{align}
Before the row error is bounded, the decomposition from the proof of
Proposition~\ref{prop:local-simplex-criterion} also retains
\begin{equation}
\label{eq:weighted-potts-conditional-entropy}
 -yK(r\Vert \mathbf u_{n-1})+\frac{y^2}{2}|v_r|^2
 \le-\{1-o(1)\}yK(r\Vert \mathbf u_{n-1}),
\end{equation}
uniformly for \(y\le L_n^{-2}\). Indeed,
\[
 |v_r|^2
 \le\frac{\|G_n\|_{\rm op}}{\log(n-1)}
       K(r\Vert \mathbf u_{n-1})
 =O\{K(r\Vert \mathbf u_{n-1})\}.
\]
We use a sharpened form of the last term. Pinsker's inequality and the
definition of \(C_n\) give
\[
 |q_i\cdot v_r|
 \le C_n\|r-\mathbf u_{n-1}\|_1
 \le C_n\{2K(r\Vert \mathbf u_{n-1})\}^{1/2}.
\]
The spectral estimate in
Proposition~\ref{prop:local-simplex-criterion} retains a fixed fraction of
the conditional entropy \(yK(r\Vert \mathbf u_{n-1})\). Young's inequality
therefore absorbs the cross term \(yc_yq_i\cdot v_r\), at the cost of
\(O(yC_n^2)\). Thus the right side of
\eqref{eq:weighted-potts-localization} may be replaced by
\begin{equation}
\label{eq:weighted-potts-localization-sharp}
 \delta_nu_i c_y^2
 -\{k_n(y)-R_n^\star c_y^2\}
 -c\,yK(r\Vert \mathbf u_{n-1})
 +O(yC_n^2).
\end{equation}
The scalar function in braces is nonnegative and, on
\(0\le y\le L_n^{-2}\), vanishes only at \(y=1/n\); its second
derivative there is \(n+O(L_n)\).  More generally, its explicit formula
gives, on this interval,
\begin{equation}
\label{eq:potts-scalar-local-stability}
 k_n(y)-R_n^\star c_y^2
 \ge
 \frac{c}{n}\{z\log z-z+1\},
 \qquad z=ny.
\end{equation}
This follows by differentiating the two sides on either side of \(z=1\);
both vanish with zero derivative at \(z=1\), while the second derivative
of the left side after the change of variables is
\(n^{-1}z^{-1}\{1+o(1)\}\).
Combining this estimate with \eqref{eq:potts-fitted-weight-profile} gives
the following dichotomy.  If \(i\notin A\), then \(u_i\le-\eta\), and the
weighted free energy is at most \(-c\eta\delta_n\) for all large \(n\).
If \(i\in A\), the active fitted weight cancels the radial term: uniformly
on the present interval,
\begin{equation}
\label{eq:potts-active-radial-cancellation}
 \delta_nu_i c_y^2+(1-y)\log c_i
 \le O_{C,J}(\delta_n/n)+o(\delta_n),
\end{equation}
because
\(\sup_{0\le y\le L_n^{-2}}\{c_y^2-\chi_n(1-y)\}
=1/(n-1)\).
Let \(\phi(z)=z\log z-z+1\).  Outside \(z=ny\in[1/2,2]\), one has
\(\phi(z)\ge c(1+z)\), so the term \(yC_n^2=zC_n^2/n\) is absorbed by
\eqref{eq:potts-scalar-local-stability}.  Inside that interval it is
\(O(C_n^2/n)=o(\delta_n)\).  Nonnegative weighted free energy therefore
implies
\[
 \phi(ny)+nyK(r\Vert\mathbf u_{n-1})=o(n\delta_n)=o(L_n).
\]
It follows that
\[
 y=o\left(\frac{L_n}{n\log L_n}\right),
 \qquad
 yK(r\Vert\mathbf u_{n-1})=o(\delta_n).
\]
Hence
\[
 y|v_r|
 \le C\{y^2K(r\Vert\mathbf u_{n-1})\}^{1/2}
 =o\{y\delta_n\}^{1/2}=o(1).
\]
Since
\(|c_y-\chi_n|=|1-ny|/(n-1)\), we obtain
\[
 \mu(p)=c_yq_i+yv_r=\chi_nq_i+o(1).
\]
Substituting \eqref{eq:potts-active-radial-cancellation} and the bounds
\(yK(r\Vert\mathbf u_{n-1})=o(\delta_n)\) and
\(y|v_r|=o(1)\) into the definition of \(F_{\log c}\) gives
\(F_{\log c}(p)=o(\delta_n)\).  Equation
\eqref{eq:weighted-certificate-square-identity} then shows that every
certificate point with value at least one lies within \(o(1)\) of
\(\chi_nq_i\).

It remains to prove uniqueness inside the localized patches.  Fix a small
constant \(r_0>0\).  In \(B(\tau,r_0)\), the Gibbs probabilities are at
most \(n^{-1+o(1)}\), and hence
\[
 \|\operatorname{Cov}_{p(w)}(q)\|_{\rm op}
 \le n^{-1+o(1)}\|G_n\|_{\rm op}=o(1).
\]
In \(B(w_i,r_0)\), the total Gibbs mass outside coordinate \(i\) is
\(n^{-1+o(1)}\), so the covariance operator is again \(o(1)\).  Therefore
\[
 \nabla^2\log\Gamma_{\widetilde\Pi}(w)
 =\operatorname{Cov}_{p(w)}(q)-I\preceq-\frac12I
\]
on every localized patch.  Each patch contains exactly one stationary
point, namely the contact already constructed there, and that point is its
unique maximum.  The diffuse and near-vertex localization arguments in this
proof now give
\(\Gamma_{\widetilde\Pi}\le1\), with equality exactly at the central and
active contacts.

The constructed measure is an NPMLE\@. Its fitted vector is unique by
strict concavity in the fitted likelihood vector.  The likelihood vectors
at its contacts are linearly independent.  Indeed, on rows indexed by
\(A\), the active-contact block has diagonal entries
\((n-1)\{1+o(1)\}\) and off-diagonal entries
\((n-1)^{-1}\{1+o(1)\}\), while the central column is \(1+o(1)\).
Choose any inactive row, which exists because \(|A|\le J<n\) for large
\(n\).  On that row the central entry is \(1+o(1)\), and every active
entry is \((n-1)^{-1}\{1+o(1)\}\).  After scaling the active columns by
\(n^{-1}\), this \((1+|A|)\)-square minor converges to an invertible
triangular matrix.  The KKT conditions for the unique fitted likelihood
vector restrict the support of every maximizing measure to these contacts;
linear independence then makes its representing weights unique.  Thus every
maximizing measure uses the same positive weights on precisely these
contacts.  Equations
\eqref{eq:buffered-simplex-count} and \eqref{eq:buffered-simplex-mass}
follow from \eqref{eq:terminal-quantitative-mass}.
\end{proof}

For the random Gaussian cloud, all assumptions of
Proposition~\ref{prop:buffered-simplex-support} hold whenever
\(d_nL_n/n^2\) is bounded below.

\begin{corollary}
\label{cor:random-buffered-simplex-support}
Suppose
\[
 d_nL_n/n^2\ge c>0,
 \qquad s_n^\star=O(1).
\]
For every fixed \(C,J<\infty\) and \(\eta>0\),
\begin{equation}
\label{eq:random-buffered-simplex-support}
 \Pbb\left\{
  \mathcal B_n(C,J,\eta),\quad
  \widehat K_n\ne1+\#\{i:\xi_{n,i}^\star>0\}
 \right\}\longrightarrow0.
\end{equation}
On \(\mathcal B_n(C,J,\eta)\), the NPMLE is also unique apart from an event
whose probability tends to zero.
\end{corollary}

\begin{proof}
Here \(L_n=o(d_n)\), and Lemma~\ref{lem:random-off-diagonal} gives
\begin{align*}
 \delta_n^\star&=\sqrt{L_n/d_n}\{1+o(1)\},&
 \mu_n^{\rm off}:=\max_{i\ne j}|q_i\cdot q_j|
   &=o_{\Pbb}(L_n^{-1}),\\
 \|G_n\|_{\rm op}&=O_{\Pbb}(L_n),&
 C_n&=O_{\Pbb}(L_n^{3/2}/\sqrt{d_n}).
\end{align*}
Consequently,
\[
 \frac{C_n}{\delta_n^\star L_n^2}=O_{\Pbb}(L_n^{-1}),
 \qquad
 \frac{C_n^2}{n\delta_n^\star}
 =O_{\Pbb}(\delta_n^\star L_n^2/n)=o_{\Pbb}(1),
\]
and \(\delta_n^\star=O(L_n/n)\to0\).  Proposition
\ref{prop:buffered-simplex-support} applies on an event whose probability
tends to one.
\end{proof}
